% Generated draft: requires mathematical and source review.
\documentclass[11pt]{article}
\usepackage[margin=25mm]{geometry}
\usepackage{fontspec}
\setmainfont{Latin Modern Roman}
\usepackage{amsmath,amssymb,mathtools,unicode-math}
\setmathfont{Latin Modern Math}
\usepackage{enumitem,needspace,xurl,hyperref}
\setlength{\parindent}{0pt}
\setlength{\parskip}{5pt}
\setlength{\emergencystretch}{3em}
\newcommand{\N}{\mathbb{N}}
\newcommand{\Z}{\mathbb{Z}}
\newcommand{\Q}{\mathbb{Q}}
\newcommand{\R}{\mathbb{R}}
\newcommand{\C}{\mathbb{C}}
\newcommand{\F}{\mathbb{F}}
\newcommand{\A}{\mathbb{A}}
\newcommand{\PP}{\mathbb{P}}
\newcommand{\E}{\mathbb{E}}
\newcommand{\eps}{\varepsilon}
\newcommand{\dd}{\,\mathrm{d}}
\begin{document}
\section*{UnsolvedMath ID: 8800082}
\textbf{Source problem code:} catalogue entry
\subsection{Definitions and mathematical statement}
For a cyclic subgroup \(\langle g\rangle\) of order \(r\), discrete logarithm inversion takes \(h\in\langle g\rangle\) and seeks \(x\) with \(g^x=h\), determined modulo \(r\). The order need not be publicly known unless the construction requires it.

Rivest asks for practical groups in which discrete logarithms have an effective trapdoor. In a customary cryptographic formulation, an efficient generation algorithm produces public group data \(\mathrm{pp}\), including a representation of \(\langle g\rangle\), and a secret \(\tau\). Public data permit efficient group operations and exponentiation; with \(\tau\), a specified inversion algorithm efficiently recovers a logarithm from \(h\). Without \(\tau\), inversion should remain computationally hard under a precise instance distribution and security assumption.

Both ingredients matter: a group in which everybody can efficiently compute logarithms has no trapdoor advantage, while a hidden description with no effective inversion is not the requested trapdoor. The source does not specify the group class, public knowledge of the order, classical or quantum model, correctness scope, distribution of logarithm instances, security definition, or concrete practicality target. The construction must therefore not be replaced by an unqualified claim about all generic or black-box groups.

The summary reports no known practical construction and mentions generic-group limitations. Such model-dependent limitations do not by themselves constitute a verified impossibility theorem for the broad source request, and no negative answer is asserted here.
\subsection{Short English statement}
Construct an efficiently usable group family whose public discrete logarithms remain hard but become efficiently invertible with secret trapdoor data. The group representation, security model, and practical efficiency criteria are unspecified.
\subsection{Sources}
\begin{itemize}
\item Local UnsolvedMath registry, permanent ID 8800082, source code catalogue entry; attributed to Ronald Rivest; no proposal year supplied; update 2026-08-06.
\item IPAM, \emph{Workshop on Number Theory and Cryptography --- Open Problems} (2006), attachment 6243, legacy PPT logical slide 17: \url{https://www.ipam.ucla.edu/programs/workshops/workshop-i-number-theory-and-cryptography-open-problems/?tab=schedule}. Registry access date 2026-07-29; PDF and legacy-PPT visual audit. Provenance reports 19 attachments, 17 PDFs and two legacy PPT decks, with repetitions and solved definitions excluded. Present status and redistribution terms require review.
\item The supplied summary mentions Maurer and later trapdoor-discrete-log work and generic-group limitations without exact references, model assumptions, or an unconditional impossibility result. Its open/negative assessment was not independently checked.
\end{itemize}
\end{document}
